\section{Dashboard Creation}

Dashboards transform raw telemetry into visual insights that help analysts monitor activity, validate detections, 
and identify anomalies at scale. For CTI analysts, threat hunters, and detection engineers, dashboards serve two 
primary purposes: they provide situational awareness, and they act as investigative starting points. A well‑designed 
dashboard highlights deviations from normal behavior, exposes suspicious trends, and allows analysts to pivot quickly 
into deeper searches. This chapter introduces the core principles of dashboard design and provides practical SPL and 
KQL examples that can be used to build effective visualizations in Splunk and Microsoft Sentinel.

% --------------------------------------------------------------------
% SECTION 1: Dashboard Design Philosophy
% --------------------------------------------------------------------

\subsection*{Dashboard Design Philosophy}

A good dashboard begins with a clear objective. Instead of trying to visualize everything, focus on the behaviors 
that matter most: process execution, authentication patterns, network activity, and high‑risk events. Each panel 
should answer a specific investigative question, such as:

\begin{itemize}
    \item “Which processes are most active?”
    \item “Which accounts are failing logons?”
    \item “Which hosts are making rare outbound connections?”
    \item “Which events indicate high‑risk behavior?”
\end{itemize}

By structuring dashboards around investigative questions, you ensure that each visualization provides actionable 
insight rather than noise.

% --------------------------------------------------------------------
% SECTION 2: Process Activity Panel
% --------------------------------------------------------------------

\subsection*{Process Activity Panel}

Process creation is one of the most valuable signals for both hunting and detection engineering. A dashboard panel 
that highlights the most common processes helps establish a baseline and makes anomalies easier to spot.

\textbf{SPL Example}

\begin{lstlisting}
index=sysmon EventCode=1
| stats count by Image
| sort - count
\end{lstlisting}

\textbf{KQL Equivalent}

\begin{lstlisting}
Sysmon
| where EventID == 1
| summarize Count=count() by Image
| order by Count desc
\end{lstlisting}

This panel helps analysts quickly identify unusual binaries or unexpected spikes in activity.

% --------------------------------------------------------------------
% SECTION 3: Authentication Trends Panel
% --------------------------------------------------------------------

\subsection*{Authentication Trends Panel}

Authentication logs reveal credential misuse, brute‑force attempts, and lateral movement. A dashboard panel that 
tracks failed logons over time provides early warning of suspicious behavior.

\textbf{SPL Example}

\begin{lstlisting}
index=wineventlog EventCode=4625
| timechart count by AccountName
\end{lstlisting}

\textbf{KQL Equivalent}

\begin{lstlisting}
SecurityEvent
| where EventID == 4625
| summarize Count=count() by bin(TimeGenerated, 1h), Account
| render timechart
\end{lstlisting}

This visualization helps identify targeted accounts or sudden increases in authentication failures.

% --------------------------------------------------------------------
% SECTION 4: Network Activity Panel
% --------------------------------------------------------------------

\subsection*{Network Activity Panel}

Network telemetry is essential for detecting command‑and‑control activity, lateral movement, and exfiltration. 
A dashboard panel that highlights rare outbound connections can reveal early signs of compromise.

\textbf{SPL Example}

\begin{lstlisting}
index=sysmon EventCode=3
| stats count by DestinationIp, DestinationPort
| where count < 5
\end{lstlisting}

\textbf{KQL Equivalent}

\begin{lstlisting}
Sysmon
| where EventID == 3
| summarize Count=count() by DestinationIp, DestinationPort
| where Count < 5
\end{lstlisting}

This panel surfaces low‑frequency connections that may indicate malicious activity.

% --------------------------------------------------------------------
% SECTION 5: High‑Risk Event Panel
% --------------------------------------------------------------------

\subsection*{High‑Risk Event Panel}

Some events are inherently high‑risk and deserve dedicated visibility. Examples include PowerShell execution, 
scheduled task creation, and privilege escalation attempts. A dashboard panel that highlights these events helps 
analysts quickly identify suspicious behavior.

\textbf{SPL Example}

\begin{lstlisting}
index=sysmon (EventCode=1 Image="*powershell.exe") OR EventCode=13 OR EventCode=7
| stats count by EventCode, Image, User
\end{lstlisting}

\textbf{KQL Equivalent}

\begin{lstlisting}
Sysmon
| where EventID in (1, 7, 13)
| summarize Count=count() by EventID, Image, User
\end{lstlisting}

This panel consolidates multiple high‑risk signals into a single view.

% --------------------------------------------------------------------
% SECTION 6: Dashboard Layout Principles
% --------------------------------------------------------------------

\subsection*{Dashboard Layout Principles}

Effective dashboards follow a few simple principles:

\begin{itemize}
    \item \textbf{Broad → Specific}  
    Start with high‑level trends (process volume, authentication failures) and move toward detailed insights.

    \item \textbf{Minimal clutter}  
    Use clear labels, simple charts, and avoid unnecessary visual noise.

    \item \textbf{Actionable panels}  
    Every panel should support a pivot into deeper investigation.

    \item \textbf{Feedback-driven refinement}  
    Dashboards should evolve based on analyst usage and emerging threat behavior.
\end{itemize}

A dashboard is not a static artifact—it is a living tool that improves as your environment and adversaries change.

% --------------------------------------------------------------------
% SECTION 7: Why Dashboards Matter
% --------------------------------------------------------------------

\subsection*{Why Dashboards Matter}

Dashboards provide continuous visibility into system behavior, highlight anomalies, and support both hunting and 
detection workflows. They help analysts validate detections, monitor trends, and quickly identify suspicious activity. 
Later chapters expand these concepts into full dashboard layouts and advanced visualizations tailored to specific 
adversary techniques.
